\documentclass[UTF8,11pt,fontset=none]{ctexart}
\usepackage[a4paper,margin=23mm,headheight=15pt,footskip=11mm]{geometry}
\usepackage{fontspec}
\IfFontExistsTF{TeX Gyre Termes}{\setmainfont{TeX Gyre Termes}}{\setmainfont{lmroman10-regular.otf}[BoldFont=lmroman10-bold.otf,ItalicFont=lmroman10-italic.otf,BoldItalicFont=lmroman10-bolditalic.otf]}
\IfFontExistsTF{TeX Gyre Heros}{\setsansfont{TeX Gyre Heros}}{\setsansfont{lmsans10-regular.otf}[BoldFont=lmsans10-bold.otf,ItalicFont=lmsans10-oblique.otf]}
\IfFontExistsTF{DejaVu Sans Mono}{\setmonofont{DejaVu Sans Mono}[Scale=MatchLowercase]}{\setmonofont{lmmono10-regular.otf}}
\IfFontExistsTF{Noto Serif CJK SC}{\setCJKmainfont{Noto Serif CJK SC}[AutoFakeSlant=0.15]}{\setCJKmainfont{FandolSong-Regular.otf}[BoldFont=FandolSong-Bold.otf,ItalicFont=FandolKai-Regular.otf]}
\IfFontExistsTF{Noto Sans CJK SC}{\setCJKsansfont{Noto Sans CJK SC}}{\setCJKsansfont{FandolHei-Regular.otf}}
\IfFontExistsTF{Noto Sans Mono CJK SC}{\setCJKmonofont{Noto Sans Mono CJK SC}}{\setCJKmonofont{FandolFang-Regular.otf}}
\usepackage{amsmath,amssymb,amsthm,mathtools,mathrsfs,esint,bm}
\usepackage{booktabs,longtable,array,calc}
\usepackage{enumitem}
\setlist{nosep,leftmargin=1.7em,topsep=3pt}
\usepackage{fvextra}
\DefineVerbatimEnvironment{Highlighting}{Verbatim}{breaklines,breakanywhere,fontsize=\small}
\usepackage{xurl}
\usepackage[unicode,hidelinks,bookmarksnumbered]{hyperref}
\usepackage{bookmark}
\usepackage{fancyhdr}
\pagestyle{fancy}\fancyhf{}
\fancyhead[L]{\small 第一科学主线：近期研究全稿}
\fancyhead[R]{\small 2026-10-07}\fancyfoot[C]{\thepage}
\renewcommand{\headrulewidth}{0.3pt}
\providecommand{\tightlist}{\setlength{\itemsep}{1pt}\setlength{\parskip}{0pt}}
\providecommand{\passthrough}[1]{#1}
\newcommand{\Ric}{\operatorname{Ric}}
\newcommand{\Vol}{\operatorname{Vol}}
\newcommand{\Rm}{\operatorname{Rm}}
\newcommand{\supp}{\operatorname{supp}}
\newcommand{\Var}{\operatorname{Var}}
\newcommand{\tr}{\operatorname{tr}}
\newcommand{\avint}{\fint}
\newcommand{\Jcal}{\mathcal J}
\newcommand{\Ecal}{\mathcal E}
\newcommand{\dd}{\,\mathrm d}
\newcommand{\Rea}{\operatorname{Re}}
\setlength{\parindent}{2em}
\setlength{\parskip}{3pt plus 1pt}
\setlength{\emergencystretch}{3em}
\linespread{1.10}\allowdisplaybreaks[2]
\setlength{\abovedisplayskip}{8pt plus 2pt minus 2pt}
\setlength{\belowdisplayskip}{8pt plus 2pt minus 2pt}
\renewcommand{\arraystretch}{1.12}
\setcounter{tocdepth}{1}\setcounter{secnumdepth}{2}
\pdfstringdefDisableCommands{\def\fint{int}\def\mathcal#1{#1}\def\mathfrak#1{#1}\def\mathscr#1{#1}\def\operatorname#1{#1}\def\varepsilon{epsilon}\def\chi{chi}\def\beta{beta}\def\rho{rho}\def\tau{tau}\def\Phi{Phi}\def\sigma{sigma}\def\kappa{kappa}\def\nabla{D}\def\mathbb#1{#1}\def\({}\def\){}\def\textsuperscript#1{#1}\def\textsubscript#1{#1}}
\title{外部研究的集中流、椭圆约化与接入条件}
\author{第一科学主线·近期研究分卷}
\date{2026年10月7日}
\begin{document}
\maketitle
\noindent\textbf{原稿编号：}R06。

本章汇总此前实际阅读的外部稿件及方法接入，沿用原记录固定的 openai/math 源版本 \texttt{adc7f1241b42e322a6451854ab7e4b4c146bf78a}。外部命题按来源陈述列出；已经用于具体推论的 Ricci 流延拓，与尚待建立的 O10/O11 接口分开。\cite{R06,R07}

\section{第342族：Taming implies compatibility on four-manifolds}\label{ch6:ux7b2c342ux65cftaming-implies-compatibility-on-four-manifolds}

论文日期：2026-09-23。\cite{OAI342}

\subsection{原稿的定理与真正可迁移的机制}\label{ch6:ux539fux7a3fux7684ux5b9aux7406ux4e0eux771fux6b63ux53efux8fc1ux79fbux7684ux673aux5236}

原稿主定理：闭连通光滑四流形上的光滑 almost-complex structure J，若被某个 symplectic form tame，则存在与同一个 J compatible 的 symplectic form。一般不保证保持原来的上同调类；\texttt{b\_2\^{}+=1} 的类保持结论是另一个推论。

证明并非直接将二形式作 J-不变投影，因为那样会破坏闭性。其反证方法构造正流 P 和 J-反不变修正 Q，使 T=P+Q 满足闭性或相应消去关系。正流带有复直线方向的测度表示，因而比一个标量能量测度保留更多几何信息。

原稿得到：

\[
\mu(B_r(x))\le Cr^2,
\]

以及近邻复直线的角度二阶矩

\[
\mathcal A_r=\iint_{d(x,y)<r}|\ell-\tau_{yx}m|^2\,d\lambda(x,\ell)d\lambda(y,m)
=O(r^4).
\]

令 E 为正二维密度集合，核心分裂定理证明 \textbf{限制流 P\_E 仍然闭合}。决定性边界项不是分别要求两个误差有很快的收敛率，而是把角度误差与横向偏离保留成乘积：

\[
Cr^{-3}\sqrt{\mathcal A_r}\sqrt{\mathcal T_r}
=r^{-3}O(r^2)o(r)=o(1).
\]

再分开正密度与零密度部分，消除修正，完成正流与 tame 形式的矛盾。

\subsection{与 O10 的具体接口}\label{ch6:ux4e0e-o10-ux7684ux5177ux4f53ux63a5ux53e3}

项目已经知道标量幅度矩关系允许径向集中。该论文提供另一种思路：\textbf{不先要求集中消失，而是证明集中携带一个闭合、定向、可配对的几何对象，再用拓扑关系限制它。}

应先检查能否从原实际曲率与径向结构构造有几何方向的二维流，证明相应闭性、正性、质量增长及角度估计。不能把 \(|\nabla v|^2dV\) 这样的标量测度直接叫成正闭二维流。

两个必须处理的差别：

\begin{itemize}
\tightlist
\item
  原稿固定一个闭流形及光滑 J-Hermitian 度量；热核、小球坐标、投影交换子常数可依赖这些固定数据。O10 的流形、J 和几何都可退化，不能直接把原常数当成密度无关常数。
\item
  O10 的几乎复结构来自可能有零点的实际旋量。必须保留零点与中间幅度层，不能由 L2 近似得到处处的全局 tame 结构。
\end{itemize}

因此优先移植的是分裂定理及其临界边界乘积，而不是主定理的名称。

\section{第351族：退化 Ricci-flat 树上的路径选择与椭圆估计}\label{ch6:ux7b2c351ux65cfux9000ux5316-ricci-flat-ux6811ux4e0aux7684ux8defux5f84ux9009ux62e9ux4e0eux692dux5706ux4f30ux8ba1}

论文：\emph{Path selection and an elliptic inequality on degenerating Ricci-flat trees}，2026-09-24。\cite{OAITree}

\subsection{原稿精确控制的对象}\label{ch6:ux539fux7a3fux7cbeux786eux63a7ux5236ux7684ux5bf9ux8c61}

固定一棵有限 Ricci-flat 四维几何树，拼接端的群非平凡，固定端部识别，根部 ALE 衰减阶 q\textgreater1。假设给定顶点上的光滑收敛、至20阶的端部度量界，以及至10阶的加权、尺度化 Ricci 小误差 M。

控制的是重整化 Einstein--Hilbert 泛函

\[
\mathscr H(g)=\lim_{R\to\infty}
\left[\int_{r<R}s_g\,dV_g-
\int_{r=R}(\partial_jg_{ij}-\partial_ig_{jj})n_i\,dS\right],
\]

结论为 \(\mathscr H(g_j)=o(M_j)\)，只对固定树及其数据成立。该泛函的长度量纲为2；O11 的 P 是四维尺度不变量。不要使用同一个字母 E 掩盖这个区别。

\subsection{最值得复用的两个步骤}\label{ch6:ux6700ux503cux5f97ux590dux7528ux7684ux4e24ux4e2aux6b65ux9aa4}

\textbf{保留核与余核的有限维约化。} 核由有限个核心测试控制；余核以紧支撑源保留；颈部临界球面模态作为独立匹配数据。不是先假设 Lichnerowicz 算子可逆，也不预设所有无穷小 Ricci-flat 变形都可积。构造实际的有限维模障碍族 h\_*(L,z)，并控制原度量到该族的距离及泛函误差。

\textbf{多尺度路径选择。} 颈长可以以不同速率趋无穷，幂---对数渐近展开不被冒充为收敛幂级数。选取保留符号和极限的解析路径，使

\[
\int\left(\sum|dz_\alpha|+\sum e^{-cL_e}|dL_e|\right)<\infty.
\]

若 F 为约化观测量，J 为约化 Ricci 平方误差，论文证明

\[
|dF|\le C\sqrt J\left(\sum|dz|+\sum e^{-cL}|dL|\right).
\]

假设失败序列满足 \(F\ne0\)、\(F\to0\)、\(F^2\ge cJ\)，路径选择保留这些关系，使 \texttt{log\textbar{}F\textbar{}} 的总变差有限，与 \(F\to0\) 矛盾。

\subsection{对 O11 的用法与缺口}\label{ch6:ux5bf9-o11-ux7684ux7528ux6cd5ux4e0eux7f3aux53e3}

可借用的是：不盲目从全度量序列跳到有限参数，而先保留全部慢模态和障碍，建立到真实参数族的定量比较，然后处理多尺度边界。

目前项目的近驻点选择并不自动提供论文要求的高阶加权范数、固定有限树或标准端部。原普通 AE 手术还可能出现平凡端群，不属于原文明确的拼接端类别。必须补入相应临界模态与估计，不能直接套用。

若应用于 O11，还须将比较误差控制为 \(o(Q_j)\)，并对实际曲率平方泛函证明自己的变分关系。\texttt{mathscr\ H=o(M)} 不等于 \texttt{P\ \textgreater{}=\ 8\ pi\^{}2\ Q-o(Q)}。

\subsection{对 O10 的另一项启发}\label{ch6:ux5bf9-o10-ux7684ux53e6ux4e00ux9879ux542fux53d1}

同族四维 Ricci 流论文将最外层曲率尺度作为主要对象，保留所有更小子尺度；最终得到

\[
cK^2\le o(1)(K+1)^2,
\]

从而 K-\textgreater0，而不是预设 K 有界后改善误差增长率。这是一种值得复用的闭合形态。但是它的时间变化、Ricci 能量和选尺度来自实际 Ricci 流，不能对静态 O10 直接使用。

\section{第348族：L2 Einstein gap 的相邻结果，不能删除曲率符号}\label{ch6:ux7b2c348ux65cfl2-einstein-gap-ux7684ux76f8ux90bbux7ed3ux679cux4e0dux80fdux5220ux9664ux66f2ux7387ux7b26ux53f7}

论文：\emph{An L2 Einstein Gap for Nonnegatively Curved Four-Manifolds}，2026-10-05。\cite{OAI348}

原稿假设闭、单连通、四维、\textbf{截面曲率非负}；若无迹 Ricci 的未体积归一化 L2 平方积分足够小，推出底层光滑流形为 S4、CP2 或 S2xS2。它使用伴随论文对精确非负截面曲率正 Einstein 度量的分类。

可借鉴的分析是高曲率尾部的动态阈值：

\[
w(t)=(|Rm|-K(t))_+,\quad K'=C_qK^2,
\qquad \int w(t)^2\le e^{Ct}\int w(0)^2.
\]

但启动它之前，原稿使用 sec\textgreater=0 得到 \texttt{\textbar{}Rm\textbar{}\ \textless{}=\ C\ s}，用 Ric\textgreater=0 的 almost-Schur 和几何比较获得体积、直径、Sobolev 控制。O10 的 s=-4 和不受控制的 W- 并不满足这些条件。O11 中总曲率常为 O(A)、局部修正总量为 O(Q)，体积平均小也不能直接替换原文的绝对小尾能量。

因此这篇是新近邻和方法来源，不是已经解决项目负 Einstein 分支的同一个定理。它也说明``无需预设几何紧性''可能通过结构条件先推出紧性；项目的实际场路线不能与此混同。

已读源码： \path{preprints/An-L2-Einstein-Gap-for-Nonnegatively-Curved-Four-Manifolds-October-5-2026/build/source/main.tex}

读取范围包含定理、初始归一化、背景高尾、Sobolev与尾传播部分。\path{lean/docs/348.md} 当前明确描述的是严格正截面曲率 Einstein 分类，不是这篇后来新增的全部 L2 gap 分析；同一族的 Lean 标记不能自动扩展到所有伴随稿件。

\section{已筛过但不宜优先套用的条目}\label{ch6:ux5df2ux7b5bux8fc7ux4f46ux4e0dux5b9cux4f18ux5148ux5957ux7528ux7684ux6761ux76ee}

\begin{itemize}
\tightlist
\item
  \textbf{341 hypersymplectic deformation}：需要三条全局闭二形式，每个非零线性组合都保持正辛性。多个 SW 解不等于这样的三元组；一条有零点的形式也不满足它。其类保持超Kähler实现很强，但目前不是 O10/O11 的直接输入。
\item
  \textbf{335 Gromov integral scalar inequality}：在四维给出标量平方对实单纯体积的下界。它不是 P 的锐利手术成本；也不允许相减两个独立下确界。
\item
  \textbf{336 spectral scalar width}：需要 \texttt{-4\ Delta+s\ \textgreater{}=\ 1}。在 s=-4 的闭代表上测试常数函数已经失败；第五线丛的正势不是这个标量算子。
\item
  \textbf{337 Cartan--Hadamard/CAT(0)}：要求截面曲率上界或 CAT(0) 几何；L2 近 Einstein 不提供这些条件，带吹起的整个流形也不能被当成原复双曲主体。
\item
  \textbf{338 Yau uniformization}：是完备非紧 Kähler、严格正全纯双截曲率的复解析统一化，不是 Yau26 的纯 Ricci 夹逼问题。
\item
  \textbf{346/355 varifold/MCF}：需要相应驻定、整性、演化方程；一般曲率缺陷测度不能直接按此调用。
\item
  \textbf{356/357 RCD}：需要指定测度、合成 Ricci 下界，某些结论还有非塌缩；四维临界 L2 Ricci 误差不能直接充当 RCD 输入。
\end{itemize}

\section{原流与残余流的最终配对}\label{ch6:ux539fux6d41ux4e0eux6b8bux4f59ux6d41ux7684ux6700ux7ec8ux914dux5bf9}

后续阅读进一步保留第342族的最终步骤。令 \[
\mathscr P'=\mathscr P-\mathscr P_E,\qquad
\mathscr T'=\mathscr T-\mathscr P_E=\mathscr P'+\mathscr Q.
\] 第二个正因子位于零密度部分，使正则化的负尾误差趋零。源文展开 \[
\begin{aligned}
0={}&\langle S_r\mathscr P,*S_r\mathscr P'\rangle
+\langle S_r\mathscr P,*S_r\mathscr Q\rangle\\
&+\langle S_r\mathscr Q,*S_r\mathscr P'\rangle
+\|S_r\mathscr Q\|_2^2.
\end{aligned}
\] 第一项下极限非负，混合项趋零，从而 \(\mathscr Q=0\)；随后使用驯服形式的正配对。这个方法保留``原对象--残余对象''的不同角色，并非要求所有集中先消失。上述是外部证明的接入记录，所需实际正流及其配对尚需在 O10 中建立。\cite{OAI342} \# 为什么第342族主定理不能直接替代 O10 的最后一步

外部定理要求一个全局光滑 J 和处处驯服它的闭二形式。它得到同一个 J 的相容辛形式，但一般不保证类不变，也不估计新度量的 P 能量。

在四维，实际非零正旋量定义的 J 满足

\[
 \Lambda_J^-=\omega_J^\perp\cap\Lambda^+,
 \qquad \Lambda_J^+=\mathbb R\omega_J\oplus\Lambda^-.
\]

SW 的 \(F_A^+\) 已沿 \(\omega_J\)；所以 \(F_A\) 在 \(\Phi\ne0\) 处本来就是 J-不变的。这里尚缺的是全局零点处理、处处正性及其在退化中的定量控制，而不是再做一次 J-不变投影。这是代数检查，不是 O10 已解。

此外，原 O11 手术族的自然供给类满足

\[
 a_\epsilon^2=A-k,
 \quad T_+=A-k-4\ell,
 \quad c_2(S^+_{a_\epsilon})[M]=\ell.
\]

当 \(\ell>0\) 时，不存在这个 Spin\^{}c 正旋量丛的处处非零截面。因此不能为了调用342，未经新的论证就把该实际旋量规范化为全局单位截面。这里没有声称这个 M 不存在任何别的 almost-complex 结构。

第342族仍值得迁移的是密度分裂定理：给定已经满足要求的正二维流 P、J-反不变 L2 修正 Q、闭性、质量增长和角度控制，则正二维密度部分 \(P_E\) 闭合。

但在 O10 中，以下两项尚未建立：

\begin{enumerate}
\def\labelenumi{\arabic{enumi}.}
\tightlist
\item
  由实际集中得到具有所需方向、闭性及几何控制的二维流。标量能量测度不自动是这种流；新稿的固定几何热核常数也不自动统一。
\item
  证明其与某个闭形式的配对同时受到正下界与零值约束。一个正闭流可以代表非零同调类，闭性本身不导致矛盾。尤其第五线丛的特征类本来就是非零的，不能将它未经证明当成精确流。
\end{enumerate}

这两项是第342族方法接入 O10 尚需的数学连接。

\section{第351族树定理进入 O11 的精确位置}\label{ch6:ux7b2c351ux65cfux6811ux5b9aux7406ux8fdbux5165-o11-ux7684ux7cbeux786eux4f4dux7f6e}

源定理固定有限 Ricci-flat 树、非平凡拼接端群、高阶加权误差和端部识别。其 E 是重整化 Einstein--Hilbert 泛函，长度尺度为2；目标 P 是四维尺度不变量。

以下三个区别必须保留。

\textbf{(i) 固定树不等于树复杂度可随 j 增长。} 原文明确不提供对不同树统一的指数。若用于数量随手术规模增长的核心，必须控制总误差为 o(Q)，不能逐个说``小误差''后不加总。

\textbf{(ii) Ricci-flat 极限不等于一般 P 极小化核心。} 一个收费 AE 单元可以标量平坦而不是 Ricci-flat，具有非零 P。源定理的 Ricci 小误差不能在这样的单元上自动满足。

\textbf{(iii) 变分估计需要正确的零阶参考值。} 源 E 在精确 Ricci-flat 几何上为零，并将这一事实与导数控制、真实有限维约化、路径选择连接。对 O11，仅有 DP 近零不能确定 P 应当等于 8\(\pi\)\textsuperscript{2}Q。后者仍需原手术标记和边界能量的真实对应。

因此树定理在这里提供方法输入，而非 O11 下界。可迁移的是核---余核保留、端部模态及合法参数约化；要真正改变 O11，必须为目标曲率平方泛函建立自己的比较，并保留 Q 尺度的拓扑参考值。当前还没有完成这一项。

\begin{thebibliography}{P99}
\small
\bibitem[R06]{R06} 本会话研究原稿，\emph{外部研究的集中流、椭圆约化与接入条件}，2026-10-07。\path{OpenAI_Math_O10_O11_Research_Transfer_20261007.md}。

\bibitem[R07]{R07} 本会话研究原稿，\emph{静态曲率间隙的定量 Ricci 流后果}，2026-10-07。\path{Line1_External_Theorems_Application.md}。

\bibitem[E13]{OAI342} OpenAI, \emph{Taming implies compatibility on four-manifolds}, 2026-09-23. 固定源版本见第6章；源稿 §§6--7. \url{https://github.com/openai/math}.

\bibitem[E14]{OAITree} OpenAI, \emph{Path selection and an elliptic inequality on degenerating Ricci-flat trees}, 2026-09-24. 同一固定源版本。

\bibitem[E16]{OAI348} OpenAI, \emph{An $L^2$ Einstein Gap for Nonnegatively Curved Four-Manifolds}, 2026-10-05. 同一固定源版本。

\end{thebibliography}
\end{document}
